\section{A shared environment for research and learning}\label{sec:people}

\subsection{Human and AI contributions to one record}

A unified platform should allow researchers to use the same mathematical objects whether a contribution originates with a person, a model, a proof assistant, or a computation. A human researcher could enter a conjecture and ask for candidate routes. An AI system could identify a missing premise and request review of the proposed bridge. A collaborator could replace that bridge with a counterexample. The graph would preserve the rejected statement, the refutation, and the narrower question that follows.

Unified representation does not erase differences in responsibility. Contributions need separate attribution for problem formulation, argument drafting, implementation, checking, interpretation, and acceptance. A model's output should not acquire the author's review status simply because it appears in the author's workspace. Likewise, a human-written conjecture remains a conjecture. Mathematical admission depends on evidence rather than the identity or confidence of the proposer.

The actor performing a recorded action and the author of an imported artifact are separate attribution fields. The case-study events record the research actor explicitly, while source pins identify the exact imported artifacts. A file path, timestamp, or content hash establishes provenance rather than authorship. The record specification in \cref{app:records} keeps these fields separate.

Researchers should control whether objectives are supplied, jointly formulated, or delegated to an agent, together with admissible evidence, review requirements, and resource policies. One project may use AI only for retrieval; another may permit candidate generation and computational testing; a third may use the graph entirely through human contributions. These uses can share a data model without requiring agreement about the role of AI in mathematics. Cross-project imports should retain their original attribution and review conditions.

The persistent record also supports collaboration across time. A new researcher can inspect the remaining obligations rather than infer them from the final paragraph of a notebook. A returning researcher can see why a plausible approach was abandoned. A specialist can review a precisely scoped interface between two arguments without claiming responsibility for an entire research program. Whether this actually reduces collaboration or review costs remains an empirical question.

\subsection{Making research paths available to students}

Students often encounter polished proofs before they learn how to select a representation, identify a promising intermediate claim, or respond productively to a failed approach. Research exposition, seminars, and mentorship already address this gap. RRE proposes an additional resource: selected, evidence-linked exploration paths that can be revisited at different levels of detail.

Three kinds of record must remain distinct. A \emph{source-derived argument} reconstructs dependencies in a written proof. A \emph{retrospective research narrative} explains a possible or remembered route after the result is known. A \emph{contemporaneous exploration record} preserves questions and outcomes as the investigation proceeds. The first two can be educational, but neither establishes the actual historical path of discovery. Access to an expert's genuine research practice requires that expert's participation, consent, and annotation of the record.

A teaching interface could stop before a recorded decision and ask the student to propose the next target. It could reveal only the premises available at that time, then compare the student's proposal with the recorded continuation. It could also ask which hypothesis a counterexample violates, whether a finite experiment supports a universal claim, or how a quantitative bound affects a later theorem. Success would mean improved reasoning on new tasks, not simply the ability to repeat the displayed sequence.

The tensor-certificate example provides two concrete exercises. First, students can compare a tensor's exact rank with the ranks of its matrix flattenings, using the family $W_n$ in \cref{sec:boundary}. They can then explain why equations in tensor entries and certificates in decomposition parameters answer different questions. Second, students can start with the general balance reduction and calculate how contracting an $n$-mode target to $k$ modes changes multiplier degrees. The degree-9 certificate provides a concrete example, while \cref{cor:degree-transfer} states the reusable bound. The exercise identifies the cost of every contracted mode and explains why the normalized degree $\delta_R(T)/n$ obeys the transfer inequality. Both exercises connect a mathematical distinction to the choice of a new research target.

Educational use should include unsuccessful and unresolved branches in a curated form. Showing only the successful path can create the false impression that research decisions were inevitable. Showing every raw event can obscure the mathematics. The appropriate interface combines a readable narrative with access to the underlying evidence and labels any reconstruction. No student learning study has yet been conducted for RRE.

\subsection{Visualization as an interface to obligations}\label{sec:visualization}

Phase II introduced a multiscale viewer for collection relationships, argument maps, individual inferences, and source passages \cite{xie2}. RRE extends that design to research targets, time, and revisions. The earlier v2 companion established a before/after and contribution-view paradigm. The new certificate-study companion applies that paradigm to six immutable graph versions and to prerequisite closures for the exact theorem, the $W_n$ family, and the nilpotent-pencil branch. \Cref{sec:visual-effects} relates these views to the recorded exploration events.

A three-dimensional overview may help navigate a large collection, while a two-dimensional layout is often more readable for checking a selected inference. The coordinates are presentation choices. Neither geometric proximity nor an attractive layout establishes mathematical relevance. An inspectable view should expose stable identifiers, distinguish proposed from supported edges, show hypotheses even when cards are collapsed, and offer a textual trace for accessibility and audit.

The companion compares saved graph states, displays the jointly required premises of an inference, and opens the corresponding statement and evidence. The six revisions were published during the investigation. A revision may summarize several events; graph version and event time therefore remain separate fields. The printed figures combine a chronology of recorded developments with views of the actual graph, allowing readers to move between the history of a question and the argument that eventually answers it.

\subsection{Personal collections and future model training}

A researcher should be able to maintain a collection aligned with personal interests and accepted review standards, while importing selected records from shared collections. Federation requires stable identities, versioned imports, and explicit treatment of conflicting statements. Personal curation must not silently change the hypotheses of an imported theorem or imply that locally accepted claims have community approval.

With permission, curated records could also support training and evaluation of AI systems: premise selection, obligation generation, scope checking, error localization, or proposing a useful next target. Training examples would need to distinguish valid derivations from unsuccessful attempts and preserve the information available at the time of a decision. Evaluation splits should separate related proof families and chronological descendants to reduce leakage. Public availability of a notebook should not be treated as consent for every training use. This is a future use of the infrastructure, not an activity performed in the present study.
